\documentclass[../main.tex]{subfiles}
\begin{document}
%==========================================TIÊU ĐỀ TẬP========================================
\begin{table}[h]
    \begin{adjustwidth}{-1cm}{}
        \begin{tabular}{>{\centering\arraybackslash}p{6cm}|>{\raggedright\arraybackslash}p{14cm}}
            \multirow{5}{6cm}
            % Cột bên trái: ÉPISODE và số tập
            {\\[-40pt]
                \flaregothic\fontsize{20pt}{20pt}\selectfont \centering
                \color{eptype}HISTOIRE\color{black}\\
                \dawnland\fontsize{72pt}{72pt}\selectfont
                \color{epnum}1
                \\[18pt]
                % \flaregothic\fontsize{14pt}{14pt}\selectfont
                % \color{epattr}BIENVENUE, \uline{\color{Banpire} Mel} !
                % \flaregothic\fontsize{18pt}{18pt}\selectfont
                % \color{epattr}ÉPISODE PERDU
                \color{black}
            }
             &
            % Dòng thứ nhất: tên tập tiếng Nhật
            {\fotskip\fontsize{18pt}{18pt}\selectfont \ruby{運}{うん}\ruby{命}{めい}の\ruby{退}{たい}\ruby{職}{しょく}\ruby{届}{とど}け！}                              \\[6pt]
            % Dòng thứ hai: tên tập tiếng Anh
             & {\cafeteria\fontsize{18pt}{18pt}\selectfont The Resignation Letter of Fate}                                                           \\[4pt]
            % Dòng thứ ba: tên tập tiếng Việt
             & {\vnmsans\fontsize{18pt}{18pt}\selectfont Tờ đơn xin nghỉ định mệnh}                                                                  \\[8pt]
            % Dòng thứ tư: Nhân vật xuất hiện
             & {\caviar{\fontsize{14pt}{14pt}\selectfont Nhân vật: \emp{kuon1} \emp{achan} \emp{haato} \emp{mio} \emp{noel} \emp{towa} \emp{watame}}} \\[8pt]
        \end{tabular}
    \end{adjustwidth}
\end{table}
%===========================================NỘI DUNG==========================================
\color{black}

\begin{center}
    (Tiếp tục từ {\flaregothic ÉPISODE 133}, phần {\abv Truyện phụ}, quyển số Mười.)
\end{center}

Ngày Ba Mươi tháng Mười Một.

Thời khắc sau cùng của mùa thu đang dần trôi qua, dần nhường chỗ cho cái rét của mùa đông đang tới. Những cơn gió lớn từ ngoài vịnh nổi lên từng hồi, lùa qua các khe phố, cuốn theo những chiếc lá cuối cùng còn sót lại. Lẽ ra đây sẽ phải là thời điểm giao mùa đẹp nhất với những tán lá đã rụng, chuẩn bị cho mùa khoác áo xanh mới, bầu không khí Giáng Sinh đang kề cận tới\dots

Nhưng kể từ khi sau tai nạn ma lực bất thường xảy ra ở văn phòng \hll\ ba ngày trước, thời tiết bỗng trở nên bất ổn đến lạ thường.

Bầu trời Akiba ảo vẫn phủ một lớp mây xám đục, lạnh và nặng nề như đang chất chứa điều gì chưa được nói ra. Không chỉ những tia sét chớp dọc ngang thỉnh thoảng gầm rú trên bầu trời vốn đã xám xịt, mà đôi lúc người ta còn cảm thấy một cơn ``rít'' vô hình len qua da thịt, khiến ngay cả đài khí tượng cũng phải đưa ra cảnh báo tạm thời về ``biến động áp lực bất thường trên tầng đối lưu''.

Buổi sáng ngày hôm đó, Kuon rời khỏi nhà sớm hơn thường lệ. Không phải vì công việc đặc biệt gì hay một buổi họp khẩn cấp nào, mà chỉ đơn giản là vì cậu tưởng hôm nay được nghỉ thông đồ án.

Trong chiếc ba-lô mà cậu luôn mang theo, ngoài vài bản thảo còn dang dở và USB dữ liệu, còn có một tờ giấy đặc biệt: tờ đơn xin nghỉ việc, thứ mà A-chan nhắn cậu mang đến trong một buổi họp hành chính.

Cậu đến văn phòng lúc mới hơn tám giờ sáng. Tầng ba vẫn yên ắng. Ngoài tiếng rì rì của điều hòa và bóng điện huỳnh quang chớp nháy nhẹ, chỉ có một người duy nhất đang ngồi làm việc.

``Chào buổi sáng, Hikawa-san. Cậu cũng đến sớm nhỉ?'' nàng Đeo kính ngước nhìn cậu, nở nụ cười nồng hậu.

Kuon đặt chiếc cặp, thở dài một hơi mệt mỏi, rồi đáp lại: ``Vâng\dots\ Hôm nay trường bọn em hoãn thông đồ án nên được nghỉ. Em tranh thủ đến sớm luôn. Với cả\dots\ cũng muốn tránh đợt gió giật ngoài kia nữa.''

``Tôi biết mọi chuyện sau ngày hôm đó rồi. Murasaki-san bảo là không có vấn đề gì, và đã nhận mọi trách nhiệm sau sự cố đó rồi.''

Cậu tặc lưỡi một tiếng, nhớ lại sự vụ hôm đó, khuôn mặt không giấu nổi vẻ ngao ngán.

A-chan bật cười. Cô nàng hiểu cảm giác bất lực đó. ``Dù gì thì cậu vẫn chăm chỉ như mọi ngày nhỉ.'' Ngay trước khi cô quay về chiếc máy tính, cô bất chợt sực nhó ra một điều: ``À, cậu mang tờ đơn ấy theo chưa?''

``Có chị ạ. Đợi em chút\dots''

Cậu lấy ra tờ đơn đã in từ tối hôm trước, gấp làm ba, đưa bằng hai tay. Khi nàng Đeo kính vừa cầm lấy, cậu nghiêng đầu hỏi: ``Nhưng\dots\ đơn nghỉ việc này là của ai thế ạ? Sao lại nhờ em in?''

Cô đồng nghiệp chỉnh lại cặp kính, rồi bình thản đáp: ``À, sắp tới công ty mình có một người xin điều động đi nơi khác. Mà theo quy trình thì họ phải nộp đơn xin nghỉ ở đây trước, rồi mới viết đơn chuyển công tác.''

Kuon nhướng mày: ``Nghe rườm rà thế\dots''

``Tôi cũng thấy thế,'' cô cười nhạt, ``nhưng cấp trên yêu cầu vậy thì phải theo thôi.''

Kuon không nói gì thêm. Cậu ngồi xuống cạnh bàn làm việc thường lệ, bật máy tính lên. Nhưng chưa kịp gõ mật khẩu thì điện thoại rung lên bần bật. Tin nhắn từ một người bạn học.

``\vie{\textit{Gì nữa đây? Giờ này chắc lại định hỏi bài\dots}}''

Cậu nhấn mở tin nhắn của người ra. Từ Yamazu, một thành viên trong nhóm làm đồ án.

``\vie{Ê Kuon, ông đâu rồi?}''

``\vie{Hả?\\Tôi đang đi làm. Có việc gì à?}''

``\vie{Hôm nay thông đồ án đấy ông quên à?}''

Tin nhắn của cậu bạn như một tiếng sét đánh ngang tai (cả về nghĩa đen khi bầu trời bên ngoài vẫn còn nổi gió và nghĩa bóng). Cậu rướn mày, tỏ vẻ nghi ngờ hơn, bèn nhắn lại: ``\vie{Tôi tưởng hôm nay nghỉ mà? Cô bảo hôm qua rồi còn gì?}''

``\vie{Có,}'' Yamazu trả lời, ``\vie{Nhưng cô đổi ý sáng nay rồi. Ông chưa đọc tin mới hả?}''

``\vie{\dots Gì?}''

Cậu vội vã mở tin nhắn trong nhóm môn học của lớp cậu. Trong đó, giáo viên bộ môn của cậu có viết: ``Hôm nay chúng ta lại thông đồ án nha! Ta sẽ bắt đầu thông từ 8h00 đến 17h30.''

Ngay khi cậu đọc xong tin nhắn kia, cậu tặc lưỡi một tiếng đầy cáu kỉnh, miệng lầm bầm chửi: ``\vie{Đệch ngựa vãi lờ\dots\ Lật mặt nhanh như bánh tráng thế kia\dots}''

Tiếng chửi của cậu vô tình lọt vào tai của A-chan. Cô nàng ngoảnh đầu về phía cậu: ``Sao vậy?''

Kuon thở một hơi dài, tỏ vẻ chán chường với tình cảnh oái oăm hiện tại: ``iáo viên bọn em đổi lịch phút chót\dots\ Bảo là sáng nay phải có mặt để thông đồ án ạ.''

``Tôi tưởng lớp cậu được nghỉ mà?''

``Em cũng tưởng như thế,'' cậu nhìn vào màn hình điện thoại, vẫn không kìm được mà mỉa mai, ``nhưng mà giáo viên của em vừa nhắn tin, lại còn lúc sáu giờ sáng nữa chứ. Ông nội em mà còn sống chắc cũng không phản ứng kịp!''

Nhìn điệu bộ của cậu, A-chan không khỏi chẹp miệng mà an ủi: ``Trường cậu đúng là dở hơi thật\dots\ Thôi thì, cậu cứ đến trường đi, mấy việc ở đây cũng không gấp lắm đâu.''

Cậu Tổng khẽ cúi đầu: ``Em căm ơn chị\dots\ Vậy thì em xin phép.'' Nhanh như cắt, cậu cất máy tính, cắp chiếc bàn lên và chạy về phía cửa.

Thế nhưng, ngay trước khi cậu cầm tay nắm cửa, cậu thoáng suy nghĩ một hồi: ``\textit{Có lẽ mình nên để lại giấy nhắn cho mọi người\dots}''

Kuon lục trong ngăn bàn, lấy một tờ giấy nhớ. Cậu viết nguệch ngoạc vài dòng, gần như theo cảm tính:

\txtbx{Hôm nay anh không thể đến để gặp mặt các em được.\\Anh đã quá mệt mỏi với chuyện này rồi.\\Đừng gọi cho anh.}

Cậu đặt mảnh giấy ngay bên cạnh tờ đơn xin nghỉ, chưa kịp đọc lại thì đã vội nhét USB vào túi và đứng lên chào đồng nghiệp lần nữa: ``Vậy em xin đi trước, A-san.''

``Ừ, đi cẩn thận. Nhớ mang ô đấy.''

Kuon rời khỏi văn phòng, bóng áo khoác sẫm màu biến mất nơi hành lang cuối tầng.

A-chan ngồi lại một mình. Gió bên ngoài rít qua khung cửa kính, rung nhẹ bức rèm kéo.

Cô nhìn tờ đơn và mảnh giấy mà Kuon để lại, khẽ thở dài: ``\textit{Cậu ta lúc nào cũng một mình xoay xở\dots\ Đúng là vừa kiên cường vừa đáng thương.}''

Cô quay lại màn hình máy tính, tiếp tục công việc. Tờ đơn và mảnh giấy kia vẫn nằm đó, bằng một cách vô tình nào đó khá cạnh nhau, đến mức nếu nhìn bên ngoài, người ta sẽ nghĩ rằng một người nào đó vừa ra đi mà không để lại lời nhắn cuối cùng cho mọi người.

Và vì lẽ đó, chúng chuẩn bị \emph{gieo nên một cơn hiểu lầm} không ai ngờ tới.

\ct

Bên ngoài khung cửa kính tầng ba, thành phố dường như vẫn chưa kịp thức dậy hoàn toàn. Dãy nhà thấp phía đông Akiba ảo bị phủ trong một lớp sương xám đục, mờ mịt như lông tơ giăng mỏng. Những bảng hiệu quán cà phê đầu phố loang lổ ánh đèn đỏ, phản chiếu lên mặt đường vẫn còn đọng nước mưa từ đêm qua.

Không gian sáng nhưng không có nắng, sáng như một tấm màn bạc căng đều, tỏa ra thứ ánh sáng lạ lùng không rõ nguồn. Ở một vài nơi xa xa, người ta có thể nghe thấy âm thanh mơ hồ của còi tàu điện ngầm hoặc tiếng còi báo động gió lớn, bị gió xoáy thổi lệch hướng khiến âm vang méo mó như giọng ai đó bị bóp nghẹt.

Bầu trời không hoàn toàn tối, nhưng tầng mây thấp đến nỗi tưởng như có thể với tay ra là chạm vào được. Những tia chớp mảnh vụn thỉnh thoảng cắt ngang không trung như những vết nứt trắng trên nền đá xám, không có tiếng sấm kèm theo, nhưng mặt kính trên các toà nhà văn phòng vẫn rung lên nhè nhẹ mỗi lần như thế.

A-chan khoác áo cardigan tím mỏng, tay cầm ly giữ nhiệt trống không, bước ra khỏi văn phòng mà Kuon vừa rời đi ít phút trước. Mái tóc tím buộc thấp hơi bay trong luồng gió thổi ngang dọc hành lang, khiến cô khẽ chau mày.

``\textit{Trời hôm nay đúng là nổi gió thật\dots\ Kiểu này lại sắp thành bão tới nơi rồi\dots}''

Thang máy vừa mở ra, cô bước vào, bấm xuống tầng trệt để lấy cà phê.

``\textit{Hy vọng là Hikawa-san sẽ ổn\dots}''

Đúng lúc ấy, cửa tầng ba lại mở ra một lần nữa.

Towa là người đầu tiên bước vào, theo sau là Haato, Noel, Watame và Mio. Họ khoác thêm áo măng-tô bên ngoài đằng trước bộ trang phục thường thấy vì thời tiết lạnh bất thường. Không khí giữa nhóm vẫn rôm rả như mọi ngày, dù ngoài trời là một trận gió kỳ lạ đang gào rú qua những khe phố vắng.

Nàng Hiệp sĩ là người cất tiếng đầu tiên, vẫn với giọng hăng hái như thường lệ: ``Chào buổi sáng mọi người! Hôm nay tớ là người đến sớm nhất đó nha\~{}!''

Cô bước hẳn vào trong, gõ gót bốt kim loại nhè nhẹ lên sàn hành lang, vẻ mặt rạng rỡ như vừa giành chiến thắng trong một trò chơi nhỏ.

Nàng Song tính bật cười thành tiếng: ``Không có đâu em ơi. A-chan còn đến trước chúng ta rồi.''

Nàng Ác ma khoanh tay lại, gật đầu ra vẻ đồng tình: ``Thậm chí còn đến trước cả Kuon-senpai nữa đấy.''

Nàng Hiệp sĩ bĩu môi, phồng má lên như một đứa trẻ bị tước mất phần thưởng: ``Mồ\~{} Tớ muốn có cảm giác vinh quang một chút cũng không được à?''

Nàng Cừu bật cười khẽ, cặp sừng nâu của cô đung đưa nhẹ khi cô lắc đầu: ``Như thế cũng không ổn đâu, Noel-senpai.''

Những tiếng cười râm ran lấp đầy lối đi, khiến bầu không khí trong văn phòng bỗng trở nên ấm áp hơn so với những gì đang diễn ra ngoài trời. Trong khi mọi người đang cởi áo khoác và tháo giày, thì nàng Sói đen nhìn quanh một lượt, rồi khẽ nghiêng đầu, ánh mắt như đang tính toán điều gì đó.

``Nhắc mới nhớ\dots\ Kuon-senpai đâu rồi nhỉ?''

Câu hỏi của cô khiến cả nhóm hơi chững lại một nhịp. Haato quay sang nhìn Towa.

``Cũng phải. Anh ấy có nói là hôm nay sẽ làm việc cả ngày mà, đúng không?''

Mio nhíu mày, giọng chậm rãi.

``Anh ấy lại có việc đột xuất à?''

Nàng Ác ma đưa tay lên chống cằm, đôi mắt màu lá mạ chuyển động như đang tua lại trí nhớ. ``Nếu mà thế thì\dots\ anh ấy phải báo trước chứ.''

Họ dừng lại ở giữa hành lang, nơi lối rẽ dẫn vào không gian làm việc chính. Gió từ hệ thống thông gió vẫn rít nhè nhẹ, và ánh sáng từ dãy đèn huỳnh quang phản chiếu lên bức tường trắng mờ tạo thành những vệt ánh sáng dao động như sóng nước.

Trong lúc mọi người còn đang suy đoán xem cậu Tổng đang ở đâu, Towa khẽ lùi lại phía sau, vô tình chạm phải cạnh bàn kê sát dàn máy tính. Một xấp giấy mỏng bị hất nhẹ, trượt khỏi mép bàn và rơi xuống sàn như một cánh hoa khô bị gió cuốn. Một tờ trong đó, in bằng giấy dày hơn, rơi riêng ra và lật úp.

Haato cúi xuống, nhặt lấy: ``Ơ\dots\ là đơn xin nghỉ việc này.''

Cô nhíu mày, lật mặt trước lên rồi cầm hai tay, đưa lên ngang tầm mắt. Cả nhóm chụm đầu vào, nhìn qua tờ giấy kia. Không khác cảnh mà Mio cùng mọi người xem qua bản báo cáo (mà sau này họ nhầm thành tuyên ngôn ăn kiêng) của một ai đó cách đây chín tháng trước\footnote{Xem lại {\flaregothic SPÉCIAL \ab XVI}, quyển số Bảy.}.

``Nghe nói là sắp có người nghỉ việc đúng không?'' Mio lên tiếng, tiến lại gần hơn, tay đặt nhẹ lên vai Haato để nhìn vào tờ giấy.

Noel bước sát tới, cau mày: ``Nhưng mà ai mới được? Tên đâu?''

Mio nheo mắt dò từng dòng, rồi chỉ vào phần đầu trang: ``Ở đây có ghi: `\dots xin phép muốn được nghỉ chức vụ quản lý\dots' này.''

Cô ngẩng lên, mắt hơi co lại vì cảm thấy có gì đó không ổn. ``Nhưng công ty mình có nhiều người làm quản lý mà, có khi là quản lý tài chính, hay thiết kế\dots''

``Không đâu,'' Towa chen vào, tay gõ gõ vào phần cuối trang giấy. ``Trong này có đoạn: `chuyển quyền quản lý cho A-chan' nữa.''

Watame cũng ghé vào, giọng khẽ run: ``Tôi thấy\dots\ cái này lạ lắm.''

Nàng Sói đen quay sang: ``Gì vậy?''

``Bình thường là\dots\ mọi việc chuyển giao quản lý đều ghi rõ: `chuyển cho A-chan và Kuon-senpai' mà. Nhưng cái này\dots\ chỉ có mình A-chan thôi.''

``Có thể là vì anh ấy ôm đồm quá nhiều việc nên mọi người không muốn để anh ấy vất vả nữa?'' nàng Song tính nói, dù chính cô cũng không chắc vào điều mình vừa nói ra.

Nhưng rồi nàng Cừu lại cúi sát hơn, chỉ vào một đoạn khác dưới phần nội dung: ``Trong đây cũng có ghi thế này: `Công việc này thực sự đã tạo cho tôi quá nhiều áp lực\dots\ từ việc chúng gây phiền hà cho tôi, tới việc thậm chí nằm mơ cũng còn thấy chúng\dots\ Tôi không thể nào chịu được nữa.''

Một khoảng im lặng ngắn. Câu chữ đó nặng nề hơn họ tưởng.

``Cái này nghe\dots'' Noel ngập ngừng, ``giống y như công việc mà Kuon-senpai hay làm.''

``Đúng rồi,'' Haato gật gù. ``Từ việc quản lý, tính toán, điều phối đến việc lập kế hoạch. Chỉ có anh ấy mới làm hết những việc đó.''

``Nhưng\dots\ anh ấy chưa từng than vãn gì cả mà,'' Mio lên tiếng, hơi bối rối. ``Kuon-senpai lúc nào cũng bình thản, đâu có biểu hiện gì là muốn bỏ việc đâu.''

``Cơ mà\dots'' nàng Hiệp sĩ hơi nhíu mày, rồi giơ tờ giấy lên sát mặt hơn, chỉ vào chữ ký ở cuối trang.

``Chữ ký này\dots''

Cô ngập ngừng, nuốt nước bọt.

``\dots rất giống của anh ấy.''

Cả nhóm đổ dồn ánh mắt về phía đó. Và quả thật, trong chữ `Kuon' đó, nét chữ nghiêng, đường bút mạnh, nhấn nhá quen thuộc ấy, không lẫn vào đâu được.

Tim họ như hụt mất một nhịp. Một cảm giác bất an bắt đầu nhen nhóm trong lòng họ.

Towa siết chặt hai tay, nói như gào lên: ``Không lẽ\dots\ \textbf{anh ấy định bỏ chúng ta sao!?}''

Noel và Watame cùng lúc kêu lên, như bật ra theo bản năng: ``Lại nữa ư!?''

Mio ngạc nhiên quay đầu: ``Lại nữa\dots\ là sao?''

Nàng Cừu bối rối một lúc, rồi đáp, giọng vẫn run run. ``Tám tháng trước\dots\ anh ấy cũng giở trò này với bọn em. Em với Subaru-senpai và Lamy-chan cũng ở đó nữa.''

Noel tiếp lời ngay lập tức, như thể hai người đã sẵn chuẩn bị một kịch bản để kể lại. ``Bọn em quyết đấu bằng trò Old Maid lúc đầu cho vui, và khi mọi người cá cược nhau như thế nào, anh ấy không ngần ngại mà đặt cả chức vụ của mình lên bàn cân luôn!''\footnote{Xem lại {\flaregothic ÉPISODE 99}, quyển số Tám.}

Haato tròn mắt: ``Và chuyện đó\dots\ kết thúc sao?''

``Anh ấy suýt thua,'' Watame cười buồn, ``và Noel-senpai phải chịu hình phạt thay\dots\ Bọn em với Lamy-chan và cả Subaru-senpai nữa\dots\ ai cũng mắng anh ấy cả buổi trời.''

``Chị không ngờ anh ấy làm vậy với tụi em\dots'' nàng Sói đen khẽ thì thầm, môi mím lại.

Bất ngờ, tiếng Towa vang lên cắt ngang mọi dòng suy nghĩ. ``\textbf{Khoan đã! Nhìn này!}''

Cô chỉ vào một mảnh giấy nhỏ hơn bị gấp lại, nằm ép sát cạnh máy tính, gần như suýt bị bỏ qua. Noel vội nhặt lấy, mở ra. Mọi người cùng chụm đầu lại, ánh mắt sững lại theo từng dòng chữ.

\txtbx{Hôm nay anh không thể đến để gặp mặt các em được.\\Anh đã quá mệt mỏi với chuyện này rồi.\\Đừng gọi cho anh.}

Nàng Hiệp sĩ đọc chậm rãi. Mỗi chữ phát ra như một nhát cắt mảnh nhưng bén.

Sự im lặng bao trùm lấy cả không gian làm việc. Không còn tiếng quạt máy, tiếng chuông ngoài đường, hay bất cứ âm thanh nào chen vào lúc đó.

Một giây sau\dots

``\dots Anh ấy định ra đi thật sao?'' Noel buông tay, mảnh giấy trượt khỏi tay cô, rơi xuống bàn.

Watame siết chặt nắm tay, nước mắt đã rưng rưng nơi khoé mắt. ``Ch-Chúng ta đã làm gì sai chứ\dots?''

Towa nhìn quanh, ánh mắt hoảng loạn: ``Là lỗi của chúng ta\dots\ ư?''

``Chị nghĩ\dots'' Haato nuốt khan, `\dots anh ấy định rời đi thật rồi\dots''

``Nhưng\dots\ chúng ta còn chưa nói gì với anh ấy mà!'' Mio vội chen vào, cố làm cho mọi người hạ nhiệt, dù giọng cô đã nghẹn lại.

Không ai có thể giữ bình tĩnh được nữa. Towa vội rút điện thoại, gọi thẳng vào số của Kuon. Cô áp máy lên tai, từng hồi chuông chờ khiến tim mọi người gần như đóng băng.

Rồi, một giọng tổng đài vang lên đầu dây bên kia một cách lạnh lùng và vô cảm: ``Thuê bao quý khách vừa gọi hiện không liên lạc được\dots''

Nàng Ác ma hạ máy xuống, mắt mở to như không tin được vào tai mình: ``Anh ấy\dots\ không nhấc máy\dots''

``Thì bởi vì\dots'' nàng Cừu thều thào, ``\dots trên mảnh giấy\dots\ đã ghi rõ rồi còn gì\dots''

``Nhắn tin thử đi!'' Towa nói gần như van nài.

Haato không trả lời. Cô chỉ im lặng, đôi mắt dán vào màn hình điện thoại như đang chờ một phép màu.

``Chị\dots\ chị không nghĩ là anh ấy sẽ đọc đâu,'' cô lặng lẽ thở ra một hơi dài. ``Anh ấy giận rồi.''

Cả nhóm im lặng. Không khí trong phòng trở nên nặng nề, như thể chính những bức tường trắng xung quanh cũng đang nén lại từng giọt lo lắng của họ.

\dots

\dots

Chính nàng Song tính là người phá vỡ sự trì trệ đó. ``Nếu ta không gọi được cho anh ấy, vậy ta sẽ trực tiếp đi tìm anh ấy vậy!''

Cô đứng thẳng người dậy, tay nắm chặt điện thoại như thể đó là vũ khí cuối cùng.

``Đó cũng là ý của em\dots'' Mio gật đầu, rồi chậm rãi nói thêm, ``\dots cơ mà, ta phải tìm ở đâu bây giờ?''

Towa bước một vòng quanh phòng, tay chống cằm, rồi nói nhanh: ``Anh ấy đi đâu được nhỉ\dots''

Một câu hỏi mà không một ai biết được. Quả thật, không phải ai cũng để ý tới mọi từng nước đi của cậu Tổng.

Ngoại trừ một người.

``Chị nhớ Kuon-senpai thường đi xe tuyến J7 để đi học và đi làm đấy,'' nàng Sói đen lên tiếng, như lật lại một chi tiết trong trí nhớ.

Chi tiết tưởng chừng nhỏ nhặt ấy lại trở nên đáng giá cho tất cả mọi người, khi họ bước đầu có thể xác định được Kuon đang ở đâu.

``Đó cũng là một manh mối,'' nàng Ác ma gật đầu. ``Còn gì nữa không?''

Nàng Cừu dè dặt nói: ``Bình thường thì giờ này anh ấy sẽ ở trường\dots''

``Nhưng mà\dots'' nàng Hiệp sĩ vừa nói vừa cúi nhìn tờ đơn trên tay, ``\dots cái giấy này rõ ràng là của anh ấy mà. Vậy tức là\dots''

Họ đồng thanh, gần như tìm được kết luận: ``\textbf{Anh ấy đã đưa đơn xin nghỉ từ sáng sớm rồi!}''

Cả căn phòng vơi bớt đi sự căng thẳng phần nào, khi cả năm cô gái dần có được manh mối.

``Vậy có nghĩa là anh ấy mới rời văn phòng không lâu đâu,'' Mio lập luận, giọng đanh lại đầy hy vọng.

``Chúng ta phải đuổi theo anh ấy thôi!'' Noel thốt lên, hai tay đập xuống bàn.

``Nhưng\dots\ bằng cách nào?'' Haato cau mày, vẫn chưa nghĩ ra hướng tiếp cận cụ thể.

``Cứ suy xét theo các xe buýt tuyến J7 là được,'' Towa đề xuất, bước nhanh đến trước máy tính.

``Ý hay đấy,'' nàng Sói đen đồng tình. ``Nhưng làm sao biết được anh ấy lên xe nào?''

Sự yên lặng lại bao trùm căn phòng thêm một lần nữa, lần này ngắn hơn. Mọi ánh mắt nhìn nhau, trong đầu ai cũng lướt qua hàng chục khả năng khác nhau.

Đột ngột, Noel lẩm bẩm: ``Mỗi chuyến cách nhau khoảng ba mươi phút đấy\dots''

``Và anh ấy mới rời phòng chưa đến hai mươi phút thôi!'' Watame kêu lên, mắt sáng rực.

Cả nhóm cùng gật đầu, như vừa được tiếp thêm động lực. Haato dậm chân lên sàn: ``Vậy là chúng ta có thể đuổi theo chiếc xe vừa rời bến gần nhất đó!''

Cùng lúc này, Towa gần như reo lên khi nhìn vào bản đồ chuyến xe buýt trên màn hình máy tính: ``\textbf{Đây rồi!} Chiếc xe có đuôi số\dots\ 75! Tuyến chạy lúc 8 giờ 10 phút!''

``Hay lắm, Towa-chan!'' nàng Cừu khen, không giấu được phấn khích.

``Vậy chúng ta đuổi theo xe đó thôi!'' Noel hô to, giơ nắm tay.

``Chúng ta sẽ bắt anh ấy phải nói ra sự thật!'' Watame tiếp lời, mắt long lanh như chuẩn bị ra trận.

``Không thể để anh ấy ra đi dễ dàng được!'' Towa nghiến răng.

Cuối cùng, Mio giơ một tay lên trời, dõng dạc tuyên bố: ``\textbf{Biệt đội tìm và bắt Kuon-senpai, xuất phát nào!}''

Bốn người còn lại đồng loạt hô vang: ``\textbf{Xuất phát!}''

Cả năm người gần như lao ra khỏi phòng, chân chạy gấp gáp qua hành lang và dãy cầu thang bộ.

Ít giây sau, thang máy vang tiếng *TINH!* và A-chan quay trở lại, tay cầm ly bình cà phê vừa pha nóng. Cô bước vào phòng, ngó quanh, thấy ghế trống không và không khí vẫn còn rung động vì những bước chân vừa vội vã rời đi.

``Mấy đứa lại bày trò gì nữa không biết\dots''

Cô thở dài, đặt bình cà phê lên bàn. Nhưng rồi ánh mắt khựng lại ở chiếc bàn làm việc bên cạnh.

Tờ đơn xin nghỉ mà Kuon đưa lúc sáng vẫn còn nằm đó. Cô cầm nó lên, lật lật vài trang, đọc lại. Có điều gì đó khiến cô hơi cau mày.

``\dots À-rế?''

Ở cuối trang đầu tiên, nơi đúng ra phải là phần ký tên của người nhận chuyển giao quyền quản lý, một chỗ để dành riêng cho chữ ký của cô, thì lại có chữ ký của chính Kuon.

``Chắc cậu ấy ký nhầm\dots'' cô lẩm bẩm, rồi gấp tờ đơn lại cẩn thận, đặt lên giá tài liệu đầu bàn.

\begin{center}
    \uline{Khoảng mười phút trước đó.}
\end{center}

Kuon chạy băng qua cổng công ty trong trạng thái gần như bán chạy, vừa chạy vừa liếc đồng hồ đeo tay, dưới làn mây xám nhạt lượn lờ sát mặt đất, ánh sáng rọi xuống từ những khe nứt mây nhạt loà một cách lạnh lẽo. Trạm xe buýt nằm ngay đầu phố, và chiếc xe tuyến J7 mang biển đuôi 75 mà cậu vẫn thường đi đang lừ lừ tiến đến.

``\textbf{Chờ đã bác tài ơi!!}'' cậu hét lớn, giơ tay vẫy mạnh để ra hiệu, nhưng\dots

Chiếc xe đó gần như không có ai xuống. Tài xế liếc qua gương chiếu hậu, thấy không ai đợi ở trạm, liền cho xe rời bến ngay trước mắt cậu.

Chiếc J7 lăn bánh rời đi trong làn khói mờ, để lại Kuon đứng thở hồng hộc, hai tay chống gối, mồ hôi lấm tấm trên trán dù tiết trời đang lạnh.

``\vie{Mẹ nó! Má, giờ lại phải chờ ba mươi phút nữa à!?}'' cậu gầm lên, cáu kỉnh đá nhẹ vào cột đèn gần đó. ``\vie{Nhóm mình lại thuộc nhóm đầu tiên lên bảo vệ mới điên chứ\dots}''

Cậu tặc lưỡi một tiếng, khuôn mặt hiện rõ vẻ chán nản. ``\vie{Thôi thì\dots}'' Cậu khẽ thở dài, ngồi xuống ghế chờ, gỡ ba lô xuống và rút điện thoại.

``\textit{Được rồi\dots\ mình sẽ nhắn cho ông ấy, bảo là sẽ đến muộn. Mình sẽ bảo họ thuyết phục cô bảo vệ sau.}''

Vừa gửi tin cho Yamazu, cậu vừa ngả lưng ra sau, nhắm mắt thở ra một hơi dài, chờ đợi xe tiếp theo.

Chính lúc đó, từ khoé mắt, cậu thấy một thứ gì đó lấp lánh trên lòng đường: một đồng xu mười yên, sẫm màu và đầy bụi.

Thông thường, chẳng ai nhặt những thứ lặt vặt như vậy giữa đường. Nhưng Kuon thì khác. Cậu có thói quen giữ lại mọi đồng xu, không phải vì keo kiệt mà vì với cậu mỗi đồng tiền đều là công sức tích cóp. Với Kuon, một yên cũng là chân quý, và tính cách này được truyền lại từ mẹ cậu, một người có tính tiết kiệm, đôi khi có phần hơi cực đoan.

Cậu bước xuống lề, cúi người định nhặt lên. Nhưng vừa chạm gần thì đồng xu ấy bất ngờ lăn đi như có một thế lực nào đó đang đẩy.

``Hở? Này- ê!'' Kuon giật mình rồi chạy theo nó, lôi cả chiếc ba-lô của mình đi.

Đồng xu cứ lăn, lăn mãi qua từng ngã tư, qua các hàng ghế đá và trạm điện công cộng. Cậu đuổi theo như bị thôi miên, không hiểu vì sao mình lại cố bắt bằng được một vật vô giá trị như vậy.

Cho đến khi nó dừng lại giữa một ngã tư lớn, rất gần với trung tâm thành phố.

Cậu con trai từ từ tiến tới, cúi người nhặt nó lên, miệng lẩm bẩm: ``Bắt được mi rồi nha\dots\ Lại thêm mười yên nữa!''

Cậu cười nhẹ, định bỏ đồng xu vào ví, nhưng rồi\dots

*ẦM ẦM ẦM ẦM*

Mọi thứ bất chợt rung lên bần bật.

Không chỉ là mặt đất. Không chỉ là lòng bàn tay. Mà toàn bộ khung cảnh xung quanh bắt đầu rung lắc, méo mó, như một thước phim bị tua ngược. Đèn tín hiệu nhấp nháy loạn xạ. Âm thanh còi xe, loa đài, tiếng gió, tất cả bị bóp méo đến mức lạ lẫm.

``\textbf{Là động đất à!? Trời mả cha nó!!}''

Trước khi cậu kịp suy nghĩ thêm, một tia sét nổ tung ngay bên kia đường, chỉ cách cậu chừng vài mét. Ánh sáng tím nhấp nháy, hoà cùng một luồng gió mạnh thổi dọc theo đại lộ chính, như thể một bàn tay khổng lồ vừa tát vào không gian.

Trong cơn hỗn loạn đột ngột đó\dots\ \emph{nó} xuất hiện.

Từ giữa lòng đường,ngay nơi đồng xu rơi xuống, một lỗ xoáy mở ra, xoáy tròn từng lớp khí quyển như vặn trục một hành tinh. Nó xanh tím ánh bạc, xoay vòng với tốc độ cao, trung tâm đen đặc như mực, còn rìa ngoài lấp lánh như lưới điện sống.

Lỗ xoáy ấy giống như bức ảnh in lỗi của thực tại, với các tia sáng kéo dài như những lưỡi dao ánh sáng đâm xuyên qua vách không gian. Không khí quanh nó bị kéo tụ về tâm lỗ xoáy, khiến mọi vật xung quanh bắt đầu trôi nổi.

\img{c-7a}

Một chiếc ghế băng bị hút bay lên. Một bảng hiệu nhà hàng bị bẻ cong. Cột điện, xe đạp, và rác thải đều bị kéo về phía miệng hố như bị quấn vào một xoắn ốc không lối thoát.

``\textit{Không lẽ\dots\ cái lỗ xoáy kia\dots\ chính là--!!}''

Cậu chưa kịp nghĩ xong, thì gió ngược ập tới, hút cả người cậu về phía trước.

Kuon hoảng loạn quay đầu, cố bám vào cột biển báo gần đó. Tay cậu nắm chặt lấy khung thép lạnh. Nhưng tiếng rít gió tăng lên thành một âm thanh như tiếng la hét xuyên suốt cả chiều không gian.

Cột thép rung mạnh, rồi\dots *RẮC!* bật gốc thẳng lên.

``\textbf{Chuyện quái quỷ gì đang xảy ra vậy trời\dots!!}''

Cậu hét lên, đôi tay trượt khỏi mặt cột, cả người bị kéo tuột vào miệng lỗ xoáy. Áo khoác bị hút ngược. Chiếc ba-lô đập mạnh vào mặt cậu. Tờ bài tập toán rơi lả tả từ chiếc cặp như tuyết giữa một ngày tận thế.

Cùng lúc đó, phía sau cậu, một toà nhà văn phòng cao tám tầng bị lật nghiêng, rồi đổ sập. Kính vỡ, gạch vụn, hệ thống dây điện loé sáng rồi đứt tung, cùng nhau bị hút vào cái lỗ đen kia.

Xa hơn, một cây cầu vượt, nơi chuyến xe buýt J7 đuôi 75 vừa đi qua, bị xé gãy. Một phần cầu gãy sụp, cuốn theo chiếc xe buýt, rơi thẳng xuống dòng kênh bên dưới.

Cảnh tượng kéo dài gần một phút. Rồi như thể chưa từng có chuyện gì xảy ra, chiếc lỗ xoáy bí ẩn kia co lại như một điểm sáng, biến mất hoàn toàn, để lại mặt đường nứt nẻ, bụi mù và một khu phố bị xé toạc như vừa trải qua chiến tranh.

Ở bên trong lỗ xoáy, nơi mà Kuon vừa hút vào, những bụi mù, những vật mà chiếc lỗ mang theo, mọi thứ đều bay tứ tung, xoay mạnh như cối xay gió gặp bão.

Không còn là thành phố.

Không còn là màu sắc.

Chỉ có bóng tối và ánh sáng đan xen, như đang rơi trong một đường hầm của ánh sáng bị làm méo, uốn lượn như sóng từ trường. Không có trọng lực. Không có phương hướng. Không có âm thanh, chỉ có cảm giác bị lật tung ra từng lớp suy nghĩ.

Xung quanh Kuon là những mảnh rời của thế giới: một tờ giấy in chữ nhật ngược, một chiếc bảng quảng cáo vỡ làm đôi, một con mèo bông bay lơ lửng.

``\textbf{Cái quái gì đang xảy ra thế\dots!!}''

Cậu hét lên, nhưng tiếng hét như bị rút lại vào cổ họng. Âm thanh vỡ ra thành những làn sóng ánh sáng, rồi tan biến giữa không trung. Cơ thể cậu rơi mãi, rơi mãi như đang chìm vào chính suy nghĩ của mình, cho đến khi mọi thứ dần tắt lịm.

\il

Trên bầu trời của một buổi sáng bất thường, mọi thứ dường như đã quay lại trạng thái ``bình thường''.

Chỉ có điều, một người đã biến mất, cùng với một phần nhỏ của thành phố Akiba ảo, không để lại bất kỳ dấu vết nào.

Và điều tồi tệ nhất là\dots

\dots không ai biết rằng chuyện đó vừa xảy ra như thế nào.

Nó chỉ xuất hiện bất thình lình tại đó, trong môt ngày xám xịt, sấm chớp đánh xuống đùng đùng, gió nổi lên mạnh như cơn bão.

Những gì nó để lại sau đó, như chúng ta thấy là một mớ đổ nát hoang tàn: dây điện thì đứt, những tia vàng vẫn còn chập chờn; những tòa nhà ở khu trung tâm thì đổ ngang dọc, kính và bụi bay khắp nơi.

Lỗ xoáy bí ẩn kia là gì?

Tại sao nó lại ở đó?

Không một ai có lời giải đáp.

\end{document}